\SourceLeaf{158}Bestimmung der Polhöhe aus $p$ untern und $q$ obern Culminationen das
Gewicht $\dfrac{4pq}{p+q}$ beigelegt wird, welches auch ganz richtig ist;
aber die Polhöhe ist nicht die Summe, sondern nur die \textls{halbe} Summe der
Zenith-Distanzen in oberer und unterer Culmination. Der Fall ist also ganz
verschieden. Allgemein lässt sich zeigen, \textls{ohne viele Mühe} aus den
Grundsätzen meiner Theoria, und \textls{äusserst leicht} bei meiner neuen
Begründung, dass, --- wenn eine Grösse $w$ aus der Combination der Grössen
$x,y,z$ etc. entsteht, so dass $w=\alpha x+\beta y+\gamma z+\text{etc.}$;
wenn dann ferner die Werthe von $x,y,z$ etc. unabhängig von einander so
ausgemittelt sind, dass diesen Bestimmungen resp. die Gewichte $p,q,r$ etc.
beizulegen sind, --- dass dann der hieraus hervorgehende Werth von $w$
eine solche Zuverlässigkeit habe, deren Gewicht durch $\frac{1}{\dfrac{\alpha\alpha}{p}+\dfrac{\beta\beta}{q}+\dfrac{\gamma\gamma}{r}+\text{etc.}} $ ausgedrückt wird.

(Gewicht ist übrigens immer dem Quadrate der Genauigkeit
direct, oder dem Quadrate des sogenannten wahrscheinlichen Fehlers
umgekehrt proportional; welches aber kein Lehrsatz, sondern bloss die
\textls{Definition} des Worts Gewicht ist.)



\Needspace{6\baselineskip}
\begin{center}
\textsc{Gauss} an \textsc{Olbers}. Göttingen, 25. Februar 1825.
\end{center}

\begin{center}
. . . . . . . . . . . . . . . . .
\end{center}

P. S. Die Vertheilung der Fehler in Gruppen, um die Summe der Quadrate zu
finden, hat \textls{allemal} die Wirkung, diese zu klein, also die Genauigkeit
scheinbar grösser zu machen, als sie ist.

Es seien die Fehler
\[
\begin{array}{lll}
a,\ a',\ a''\ .\ .\ .\ .\ ., & \text{Anzahl}=\alpha, & \text{Mittel}=A, \\
b,\ b',\ b''\ .\ .\ .\ .\ ., & \text{»}\ \ \beta, & \text{»}\ \ B, \\
c,\ c',\ c''\ .\ .\ .\ .\ ., & \text{»}\ \ \gamma, & \text{»}\ \ C, \\
& \text{etc.,} &
\end{array}
\]
so ist genau
\[
aa+a'a'+a''a''\ldots=\alpha A^2+(a-A)^2+(a'-A)^2+(a''-A)^2+\text{etc.}
\]
\[
bb+b'b'+b''b''\ldots=\beta B^2+(b-B)^2+(b'-B)^2+(b''-B)^2+\text{etc.}
\]
\[
cc+c'c'+c''c''\ldots=\gamma C^2+(c-C)^2+(c'-C)^2+(c''-C)^2+\text{etc.}
\]
\[
\text{etc.}
\]

